\documentclass[12pt,a4paper,titlepage]{article}
% $Id: example-techreport.tex 1568 2012-11-20 16:01:57Z foley $
% I can comment things out so that the compiler ignores them by using
% these % symbols.  They comment everything after the symbol
% If I want to put % into my text, I need to do this: \%
% I also need to be careful about using other characters like &, which
% will cause problems.
% If there is an error, usually putting the \  in front will fix it.

\usepackage{setspace}
% Create \onehalfspacing command which will set spacing to 1.5

% the \usepackage{}  lines allow me to make use of code other people
% have written.  The ones that are most commonly found just need
% \usepackage{}.  Others may need you to install the package.
% the things in [] square brackets are options to the packages
\usepackage[utf8]{inputenc}  % this allows latex to understand UTF
                             % encoding in the file. This is
                             % important for Icelandic characters.
%\usepackage[icelandic]{babel} % this makes figures and tables have
%their Icelandic names
\usepackage[T1]{fontenc} % this enables icelandic characters in the output

\usepackage{amsmath} % these let you make neat equations.  AMS is the
                     % American Math Society
\usepackage{amsfonts}
\usepackage{amssymb}
\usepackage[amssymb]{SIunits}

% this allows us to format units correctly
% we need the amssymb option to avoid conflict on \square
%\usepackage[amssymb]{siunitx}

% this creates the Pseudocode environment for algorithms
\usepackage{pseudocode}

% this allows me to include code snippets and make them pretty
\usepackage{listings}

% this allows me to adjust the margins
\usepackage[left=2.00cm, right=2.00cm, top=2.00cm, bottom=2.00cm]{geometry}


%\usepackage[today,fancyhdr]{svninfo} % grab the SVN revision information
% This package allows you to include svn information into your document
% You will need to set the svn:keywords property to have Id and HeadURL for it to work.
%% The ``today'' option sets the current day to the latest SVN date
%
%\setlength{\headheight}{15pt} % to make fancyhdr not give a warning
%% These lines are filled in by SVN if you set the keywords
%\svnInfo $Id: example-techreport.tex 1568 2012-11-20 16:01:57Z foley $
%\svnKeyword $HeadURL: https://projects.cs.ru.is/svn/mechatronics1-2012/Templates/LaTeX-TechReport/example-techreport.tex $
%\pagestyle{fancyplain}
%% This avoids putting the information on the first page

\usepackage{graphicx}
% this package allows \includgraphics to work for putting pictures in
% the document.  The older version is graphics, you should use
% graphicx instead.

\usepackage{url}
% this package allows URLs, which show up in the bibliography and
% possibly elsewhere.  you use it like this: \url{www.google.com}

\usepackage{hyperref}
% Make magic hyperlink

\author{Páll Helgason\\Sindri Ólafsson\\Sveinn Elmar Magnússon\\Þór Tómasarson}  % My name, for the titlepage
\title{\includegraphics{graphics/ru-logo}\\\vspace{10mm}
T-411-MECH Mechatronics 1\\M\&M Color Sorting Machine\\Team Proposal}  % The title, for the titlepage


\begin{document} % this tells the compiler that it is time to make
                 % text to print instead of just getting ready.
\maketitle  % make a title page from the Title, Date, and Author
\onehalfspacing % 1.5 spacing

This report is a proposal of a mechatronic project to do as a
final project in the course T-411-MECH Mechatronics for the fall of 2014.
 

\section{Andy - Stuff}



\subsection{Mechanics} 



\subsection{Sensors}


\subsubsection{Controls} 
 

\subsection{Bill of material} % (fold)
\begin{center}
    \begin{tabular}{ | p{1.5cm} | p{1.4cm} | p{1.5cm} | l | p{5.5cm} | p{5cm} |}
    \hline
    Item & Already have & Quantity & Cost & URL & Summary \\ \hline
    Servo & \checkmark & 2x & 8.95\$ & \url{https://www.sparkfun.com/products/9065} & 180\degree servo from sparkfun. \\ \hline
    LCD Display & \checkmark & 1x & 13.95\$ & \url{https://www.sparkfun.com/products/9053} & Basic 16x2 Character LCD - Black on Green 3.3V. \\ \hline
    Color sensor & - & 1x & 25\$ & \url{https://www.sparkfun.com/products/10904} & HDJD-S822 converts light to R,G,B voltage output. From sparkfun \\ \hline
	Building material & \checkmark & - & - & - & Nylon, Plexy glass, Wood, Metal, etc. \\ \hline
    \end{tabular}
\end{center}
 




% subsection subsection_name (end)
\bibliographystyle{IEEEtran}
\bibliography{references}

\end{document} % this tells the compiler that we are done=======
